\documentclass[border=0pt]{standalone}
% The shared preamble of every docs diagram (docs/diagrams/<name>.tex).
% dvisvgm's driver first, so patterns and opacity survive into the SVG.
\def\pgfsysdriver{pgfsys-dvisvgm.def}
\usepackage[siunitx,RPvoltages]{circuitikz}
\usepackage{xcolor}
% Latin Modern: outline fonts for every glyph a label uses, µ included.
\usepackage[T1]{fontenc}
\usepackage{lmodern}
\usetikzlibrary{arrows.meta,backgrounds,calc,patterns,positioning}
\tikzset{
  % A white card, so diagrams read in mdBook's dark themes too.
  every picture/.append style={show background rectangle,
    background rectangle/.style={fill=white,inner sep=4pt}},
  % Layout geometry: layers told apart by fill and hatching, never colour alone.
  layer a/.style={draw=black,line width=0.6pt,fill=gray!20},
  layer b/.style={draw=black,line width=0.6pt,fill=gray!45},
  layer c/.style={draw=black,line width=0.6pt,fill=white,postaction={pattern=north east lines}},
  layer d/.style={draw=black,line width=0.6pt,fill=white,postaction={pattern=crosshatch dots}},
  marker/.style={draw=black,dashed,line width=0.6pt},
  % A measured distance and its label.
  dim/.style={{Latex[length=2mm]}-{Latex[length=2mm]},line width=0.5pt},
  dim label/.style={font=\small,fill=white,inner sep=1pt},
  % Where a rule fails: thick, dashed and labelled, not only red.
  violation/.style={draw=red!75!black,line width=1.4pt,dash pattern=on 3pt off 2pt},
  violation label/.style={font=\small\bfseries,text=red!75!black},
  caption/.style={font=\small},
}

\begin{document}
\begin{circuitikz}
  % Layout from above: one nfet, its poly gate across the diffusion.
  \fill[layer a] (0,0) rectangle (4,2.4);
  \fill[layer b] (1.6,-0.6) rectangle (2.4,3.0);
  \fill[layer c] (0.4,0.9) rectangle (1.0,1.5);
  \fill[layer c] (3.0,0.9) rectangle (3.6,1.5);
  \node[caption,anchor=north] at (0.7,0.8) {source};
  \node[caption,anchor=north] at (0.7,0.4) {tail};
  \node[caption,anchor=north] at (3.3,0.8) {drain};
  \node[caption,anchor=north] at (3.3,0.4) {outp};
  \node[caption,anchor=south] at (2,3.05) {poly gate, inp};
  \node[caption,anchor=east] at (-0.1,1.2) {diff};
  \draw[dim] (4.3,0) -- (4.3,2.4) node[dim label,midway,right=2pt] {$W = \SI{2}{\micro\metre}$};
  % The two capacitors pex_device adds: (0.3 + 0.2) aF/um times W on each side.
  \draw (9,2.6) node[circ,label=above:{inp}]{} to[C, l_=\SI{1}{\atto\farad}] (7.4,0.2) node[circ,label=below:{tail}]{};
  \draw (9,2.6) to[C, l=\SI{1}{\atto\farad}] (10.6,0.2) node[circ,label=below:{outp}]{};
\end{circuitikz}
\end{document}
